\section{Практичне дослідження генерації випадкових значень}


\subsection{Використання звичайного генератора псевдовипадкових чисел}

Для дослідження відтворюваності псевдовипадкових чисел було використано стандартний модуль \texttt{random} 
мови програмування Python. Спочатку було встановлено фіксоване початкове значення \texttt{seed = 42}, 
після чого згенеровано послідовність із десяти випадкових цілих чисел у діапазоні від 0 до 255. 
Далі генерацію було повторено з тим самим початковим значенням, а отримані послідовності порівняно між собою.

\begin{lstlisting}[language=Python]
import random

seed = 42
n = 10

random.seed(seed)
sequence1 = [random.randint(0, 255) for _ in range(n)]

random.seed(seed)
sequence2 = [random.randint(0, 255) for _ in range(n)]

print("I try:", sequence1)
print("II try:", sequence2)
print("Same?", sequence1 == sequence2)
\end{lstlisting}

У результаті експерименту було встановлено, що обидві послідовності повністю збігаються. 

\begin{lstlisting}[language=Python]
I try: [57, 12, 140, 125, 114, 71, 52, 44, 216, 16]
II try: [57, 12, 140, 125, 114, 71, 52, 44, 216, 16]
Same? True
\end{lstlisting}

Це пояснюється тим, що генератор \texttt{random} є детермінованим і за однакового початкового 
значення та однакових умов генерує ту саму послідовність чисел.

Якби отримані значення використовувалися для формування криптографічних ключів, 
така відтворюваність могла б становити загрозу безпеці. Якщо зловмисник знає початкове значення 
або може його вгадати, він може повторити роботу генератора та отримати ті самі значення, 
що використовувалися під час створення ключа. Це дозволило б звузити пошук можливих ключів 
і потенційно призвести до розкриття захищених даних. Тому звичайний генератор \texttt{random}, 
призначений для моделювання та статистичних задач, не слід використовувати для генерації криптографічних ключів. 
Для цього необхідно застосовувати криптографічно стійкі генератори випадкових чисел.

\subsection{Використання криптографічно стійкого генератора}

Для практичної демонстрації роботи криптографічно стійкого генератора випадкових чисел (CSPRNG)
використано стандартну бібліотеку \texttt{cryptography} у мові Python~\cite{cryptography, SP800-90A}.

Імпорт необхідних модулів:
\begin{lstlisting}[language=Python]
from cryptography.hazmat.primitives.ciphers.aead import AESGCM
\end{lstlisting}

\subsubsection{Генерація криптографічних ключів}

\begin{lstlisting}[language=Python]
key1 = AESGCM.generate_key(bit_length=256)
key2 = AESGCM.generate_key(bit_length=256)

print("Key 1 (hex):", key1.hex())
print("Key 2 (hex):", key2.hex())
\end{lstlisting}

Результат запуску:

\begin{lstlisting}[language=bash]
Key 1 (hex): a1e668462ba47835e8f5d20e6753997a9651ad0c05fd492fe37e47b671c90583
Key 2 (hex): 75d8e1eb2917f2c811df249e01661114570551273bb1851d63f7947aaa3115f1
\end{lstlisting}

\subsubsection{Повторний запуск програми}

Під час повторного запуску програми (повторного виконання скрипту) отримано такі резултати:

\begin{lstlisting}[language=bash]
Key 1 (hex): d7d04a1ff68c38309478f8221868df67fa601470edbc2950b176611de874ce25
Key 2 (hex): 85fb6c2041ab72c0c2dcaaf5b4fb05e64712585629c7b864a02755fad8d70312
\end{lstlisting}

\subsubsection{Порівняння результатів з експериментом 3.1}

У звичайному некриптографічному генераторі (розглянутому в п. 3.1) за встановленого фіксованого
початкового значення (\textit{seed}) згенеровані послідовності повністю відтворюються при кожному
перезапуску програми. Натомість при використанні CSPRNG ключі щоразу є повністю унікальними та
непередбачуваними, а їх відтворення чи повторення після перезапуску є неможливим.

\subsubsection{Джерело криптографічно стійких випадкових значень}

Програмний інтерфейс отримує випадкові значення безпосередньо з операційної системи, наприклад, через
криптографічні функції \texttt{BCryptGenRandom} у Windows або системний виклик \texttt{getrandom()} у
Linux~\cite{SP800-90A, cryptography}. Ядро ОС динамічно акумулює фоновий апаратний шум (ентропію) від
переривань пристроїв введення, таймерів і системних подій, гарантуючи високу якість та непередбачуваність
ключового матеріалу.
